\section{Faithful character tuples on a common source}\label{sec:characters}
All characters in this section are complex irreducible characters. A
faithful tuple means that the intersection of its kernels is trivial.
Slots may contain the trivial character; this permits harmless padding.
Binary-primary linear means every power-of-two order, rather than only
order two.

\subsection{Whole-socle feasibility}
Write the abelian part of $\Soc(Q)$ as a direct sum of simple
$\F_pQ$-modules over its various primes. For a nontrivial simple type
$X$ let $m_X$ be its multiplicity and
$e_X=\dim_{\End_QX}X$. For each prime let $m_p$ be the dimension of
the central elementary $p$-socle and let $f_p$ be the dimension of its
intersection with $Q'$. Put
\[
 T=\max\left(\{f_p\}_p,
       \{\lceil m_X/e_X\rceil\}_{X\ne1},
       \{1:\Soc(Q)\text{ has a nonabelian factor}\}\right),
 \qquad M_o=\max_{p\text{ odd}}m_p,
\]
with empty maxima zero.

\begin{theorem}[Three-category feasibility]\label{thm:character-feasibility}
A faithful tuple with $r$ binary-primary linear slots, $l$ unrestricted
linear slots, and $t$ general irreducible slots exists if and only if
\begin{equation}\label{eq:character-feasibility}
 t\geq T,\qquad l+t\geq M_o,\qquad r+l+t\geq m_2.
\end{equation}
The trivial group permits the empty tuple.
\end{theorem}
\begin{proof}
Every nontrivial normal subgroup intersects the socle nontrivially.
It therefore suffices to detect the entire socle with the joint tuple.
For the abelian socle, Clifford theory says that the restriction of an
irreducible contains one orbit of linear characters. In a simple
isotype $X^{m_X}$, a single orbit spans at most $e_X$ independent
multiplicity directions. To see this, identify the image algebra on
$X$ with $M_{e_X}(\End_QX)$; a vector in the isotype supplies at most
$e_X$ rows in its multiplicity space. Conversely $t$ vectors generate
that isotype whenever $m_X\leq te_X$, by choosing rows to span the
whole multiplicity space. The dual module has the same dimensions
and multiplicity, so this criterion applies equally to character orbits.

Linear characters kill $Q'$. All nontrivial abelian socle types lie
in $Q'$, because their commutator with $Q$ is the whole simple type.
Thus they require the general slots. So does each derived-central
direction, and a nonabelian socle factor requires at least one general
slot. Each irreducible is scalar on a central elementary $p$-group,
so a slot detects at most one independent central direction. A
binary-primary linear character kills every odd central element.
These facts prove necessity of \eqref{eq:character-feasibility}.

For sufficiency, choose bases of each central $p$-socle adapted to its
derived-central subspace. Assign its first $f_p$ dual coordinates to
the general slots, and as many remaining coordinates as those slots
can hold. At odd $p$, put the remaining coordinates in the $l$ linear
slots; at $p=2$, use the $r+l$ linear slots. The inequalities guarantee
room. The functionals assigned to linear slots vanish on the entire
intersection of the abelian socle with $Q'$. They extend first to its
image in $Q^{\rm ab}$ and then to complex linear characters of
$Q^{\rm ab}$, since a character of a subgroup of a finite abelian
group extends to that group. For a binary-primary slot, extend inside
the binary component and take the trivial character on all odd
components. Higher binary roots may be necessary and are allowed.

On each nontrivial simple isotype choose $t$ dual vectors whose
$Q$-orbits generate its full dual, using the matrix-algebra criterion
above. Combine their assignments over all primes with the chosen
central functionals to obtain $t$ characters of the abelian socle.
If the nonabelian socle is nontrivial, choose a nontrivial irreducible
on each of its simple direct factors and take their tensor product.
It is faithful: the simple factors are centreless, and a tensor product
can have scalar cancellation only from scalar images. Assign this
character to one general slot and use the trivial character there in
the other slots.

Each slot now specifies an irreducible of the whole socle. Induce it
to $Q$ and choose an irreducible constituent. By Clifford theory, its
restriction contains every conjugate of the specified socle character;
its kernel on the socle is precisely the corresponding normal core.
The chosen orbit spans and central coordinates make the intersection
of these cores, together with the extended linear kernels, trivial.
The resulting tuple is faithful on $Q$.
\end{proof}

For the source costs below, the cheapest feasible profile is
\begin{equation}\label{eq:cheapest-character-profile}
 t=T,\qquad l=\max(M_o-T,0),\qquad
 r=\max(m_2-\max(M_o,T),0).
\end{equation}
Indeed replacing an unrestricted linear slot by a general one increases
its coefficient from $(\log3)/3$ to $(\log5)/3$, and replacing a
binary-primary linear slot by either type also increases its coefficient.
For fixed $t$ choose the least feasible $l$ and then the least $r$;
increasing $t$ cannot improve that cost.

\subsection{Joint kernel counting}
\begin{proposition}\label{prop:character-count}
If $Q$ has a specified faithful tuple with profile $(r,l,t)$, then
for every actual $J\leq S_b$,
\begin{equation}\label{eq:character-count}
 |\Epi(J,Q)|\leq|\Aut Q|\,2^{rb/2}3^{lb/3}5^{tb/3}.
\end{equation}
\end{proposition}
\begin{proof}
For each onto kernel choose one epimorphism and inflate the fixed
tuple. Its joint kernel is that original kernel, so distinct kernels
give distinct ordered tuples of characters of the same $J$.
There are at most $2^{b/2}$ binary-primary linear characters: a Sylow
$2$-subgroup of $J$ maps onto the binary component of $J^{\rm ab}$,
and its abelianization has order at most $2^{b/2}$.
The total number of linear characters is $|J^{\rm ab}|\leq3^{b/3}$
by \cite[Theorem 2.8]{RoneyDougalTracey2025} (with the perfect case
immediate). The number of all irreducibles is
$k(J)\leq5^{b/3}$ by \cite[Theorem 1.1]{GaronziMaroti2015}, including
$b\leq3$ by direct inspection. Enlarge the joint tuple set to the
product of these three lists. Finally each onto kernel has exactly
$|\Aut Q|$ maps.
\end{proof}

\begin{proposition}[Several anchors, one preimage]\label{prop:anchor-count}
In a specified faithful tuple of profile $(r,l,t)$ choose an anchor
subtuple whose joint kernel $K$ is a $p$-group. Let $u$ be the number
of general slots among the anchors, $g=t-u$, and $d=[Q:O_p(Q)]$.
Then
\begin{equation}\label{eq:anchor-count}
 |\Epi(J,Q)|\leq|\Aut Q|d^g
        2^{rb/2}3^{lb/3}5^{ub/3}(38/25)^{gb}.
\end{equation}
If $Q$ itself is a $p$-group the empty anchor is allowed, with
$u=0,d=1$.
\end{proposition}
\begin{proof}
The inflated anchor tuple determines $M=\varphi^{-1}(K)$.
Since $K\leq O_p(Q)$ and $O_p(Q/K)=O_p(Q)/K$, it determines the
single subgroup $F=\varphi^{-1}(O_p(Q))$ as the preimage of the
characteristic $p$-core of $J/M$.
Each remaining general character, restricted to $F$, has constituents
with $p$-group image. All such constituents factor through the canonical
maximal $p$-quotient of $F$. An actual Sylow $p$-subgroup $P\leq F$
maps onto that quotient, so its irreducible list has size at most
$k(P)\leq(38/25)^b$ by the nilpotent class bound
\cite[Theorem 1.5]{Maroti2005Classes}. No faithful degree is assigned
to the abstract quotient. Frobenius reciprocity bounds the number of
irreducibles of $J$ above a fixed constituent by $[J:F]=d$:
their dimensions are at least the constituent's dimension, whereas
its induction has $d$ times that dimension.
The linear slots use the global lists of Proposition~\ref{prop:character-count},
and only the $u$ general anchors use the general list. Enlarge the
joint lists after $F$ has been determined. Faithfulness recovers the
onto kernel, and one factor $|\Aut Q|$ counts the maps.
For a $p$-group target $F=J$ is already known and no anchor is needed.
\end{proof}

\begin{corollary}[Complete finite character entries]\label{cor:character-forward}
Fix a finite menu of actual actions $U$. For each literal $N\lhd U$
choose one tuple or anchor certificate whose coefficient of $b$ in
\eqref{eq:character-count} or \eqref{eq:anchor-count} is at most
$h(w)/8-\delta$, where $w=\deg U$ and $\delta>0$ is common to the
menu. The complete accepted-orbit union has zero scalar and a
forward kernel with exponentially small row, in original weights.
\end{corollary}
\begin{proof}
Sum the chosen bounds over the actual normal axes of each $U$,
always on the same complete $J$. Their fixed target constants have
a finite sum $D_U$, so $Z_J(U)\leq D_U2^{(h(w)/8-\delta)b}$.
The pointed cold row is
$\Delta(n,b)\sum_{\deg U=n-b}D_U2^{(h(w)/8-\delta)b}/a(U)$.
By Lemma~\ref{lem:uniform-delta} its logarithm is at most
$-\delta b+O_W(\log(n+2))$, and $b\geq n-W$ for a fixed width cap
$W$. The complete finite sum is exponentially small.
\end{proof}
